\chapter{Network protocol v1}\label{cap:protocol}

A summary of the protocol; the full specification is \cmd{PROTOCOL-v1.md},
included in every release. Any client that implements it works with the
device, in any language.

\section{Transport}

HTTP/1.1 over TCP, port 8080 (8443 with TLS). Control and status in JSON;
frames in a binary stream of their own. Every path carries the \cmd{/v1}
prefix: within v1, fields and endpoints are only ever added, never removed.

\section{Endpoints}

\begin{center}
\begin{tabularx}{\linewidth}{@{}l X@{}}
\toprule
\ep{GET}{/v1/health} & available without a token; \cmd{\{"ok": true\}} \\
\ep{GET}{/v1/device} & versions (service, driver, bitstream, core), capabilities, limits, TLS \\
\ep{GET}{/v1/status} & state (\cmd{idle}, \cmd{decoding}, \cmd{error}), decode in progress, last sequence \\
\ep{POST}{/v1/decode} & the clip in the body; the response is the frame stream \\
\ep{POST}{/v1/reset} & recovery: resets the core and aborts the decode in progress \\
\ep{*}{/v1/debug/...} & diagnostics, only with \cmd{-{}-debug}, no compatibility guarantee \\
\bottomrule
\end{tabularx}
\end{center}

One decode at a time: a second \ep{POST}{/v1/decode} gets 409 while the first
one runs. The other endpoints always answer.

\section{The frame stream}

A sequence of records, integers \emph{little-endian}:

\begin{center}
\begin{tikzpicture}[x=11.5mm, font=\sffamily\scriptsize,
    f/.style={draw=#1, fill=#1!10, minimum height=9mm, anchor=west, inner sep=1pt, align=center}]
  % row 1: the 12-byte header every record starts with, one unit per byte
  \foreach \x/\w/\t in {0/4/{magic\\\texttt{"M2FR"}}, 4/1/type, 5/1/flags,
                         6/2/{header\\\_len}, 8/4/{payload\_len}}
    \node[f=m2blue, minimum width=\w*11.5mm] at (\x, 0) {\t};
  \foreach \b in {0,4,5,6,8,12} \node[lbl, below] at (\b, -0.48) {\b};
  \draw[decorate, decoration={brace, amplitude=4pt}, m2grey]
    (0, 0.55) -- node[above=4pt, lbl]{common header: 12 bytes} (12, 0.55);
  % row 2: what follows
  \node[f=m2teal, minimum width=5*11.5mm] at (0, -1.7)
    {type-specific fields\\FRAME: 20 bytes · END: none};
  \node[f=m2amber, minimum width=7*11.5mm] at (5, -1.7)
    {payload\\FRAME: I420 frame · END: JSON summary};
  \draw[decorate, decoration={brace, mirror, amplitude=4pt}, m2grey]
    (0, -2.25) -- node[below=4pt, lbl]{\texttt{header\_len} $-$ 12} (5, -2.25);
  \draw[decorate, decoration={brace, mirror, amplitude=4pt}, m2grey]
    (5, -2.25) -- node[below=4pt, lbl]{\texttt{payload\_len}} (12, -2.25);
\end{tikzpicture}
\end{center}

\begin{itemize}
  \item \textbf{FRAME} (type 1): display index, decode index, width, height,
    picture type and structure; the payload is the frame in I420.
  \item \textbf{END} (type 2): the last record, with a JSON summary: frames
    sent, whether the decode was complete, the decoder's error bits and
    per-stage timings.
\end{itemize}

A client must use \cmd{header\_len} and \cmd{payload\_len} and skip the fields
and record types it does not know: that is how the protocol grows without
breaking old clients.

\section{Error codes}

\begin{center}
\begin{tabular}{@{}c l l@{}}
\toprule
\textbf{HTTP} & \textbf{code} & \textbf{when} \\
\midrule
400 & \cmd{bad\_request} & invalid parameter, or the body is not MPEG-2 \\
401 & \cmd{unauthorized} & token missing or wrong \\
404 & \cmd{not\_found} & unknown path \\
409 & \cmd{busy} & a decode is already running \\
429 & \cmd{too\_many\_attempts} & too many wrong tokens \\
503 & \cmd{decoder\_unavailable} & the hardware does not respond \\
\bottomrule
\end{tabular}
\end{center}
